\section*{Bab \ref{ch:g8:plastics} --- Plastik dan Bahan Sintetis}

\begin{solution}{exo:g8:plastics:1}
Alami: katun, wol. Artifisial: kertas, viskosa. Sintetis: nilon,
polistirena.
\end{solution}

\begin{solution}{exo:g8:plastics:2}
Dari minyak bumi (atau gas alam).
\end{solution}

\begin{solution}{exo:g8:plastics:3}
PET (polietilena tereftalat): botol air dan minuman soda.
\end{solution}

\begin{solution}{exo:g8:plastics:4}
\omterm{def:g8:plastics:thermoplastic}{Termoplastik} melunak setiap kali dipanaskan dan dapat dicetak ulang;
\omterm{def:g8:plastics:thermoplastic}{termoset} mengeras untuk selamanya pada kali pertama dan menghangus
alih-alih melunak.
\end{solution}

\begin{solution}{exo:g8:plastics:5}
Polipropilena (PP), sebuah \omterm{def:g8:plastics:thermoplastic}{termoplastik}: wadah itu dapat dilelehkan dan
dicetak ulang.
\end{solution}

\begin{solution}{exo:g8:plastics:6}
Segitiga itu hanya menyebutkan \omterm{def:g8:plastics:plastic}{plastiknya}, untuk membantu pemilahan.
Apakah benda itu didaur ulang bergantung pada apakah benda itu
dikumpulkan dan apakah ada pabrik yang mendaur ulang \omterm{def:g8:plastics:plastic}{plastik} itu.
\end{solution}

\begin{solution}{exo:g8:plastics:7}
Wajan menjadi panas: \omterm{def:g8:plastics:thermoplastic}{termoplastik} akan melunak, \omterm{def:g8:plastics:thermoplastic}{termoset} tidak.
\end{solution}

\begin{solution}{exo:g8:plastics:8}
$\frac{353}{460} \approx 0.77$: sekitar \qty{77}{\percent}.
\end{solution}

\begin{solution}{exo:g8:plastics:9}
PVC mengandung klorin: membakarnya melepaskan hidrogen klorida, gas
yang korosif, selain jelaga dan karbon monoksida dari \omterm{def:g8:combustion-and-fuels:complete}{pembakaran tidak
sempurna} api \omterm{def:g4:what-a-fire-needs:burning}{pembakaran} sampah di halaman.
\end{solution}

\begin{solution}{exo:g8:plastics:10}
$300 \times 36 = 10\,800$ botol; $10\,800 \times 25 = 270\,000$\,g, yaitu
\qty{270}{kg}.
\end{solution}

\begin{solution}{exo:g8:plastics:11}
Karbon dioksida dan air. Karbonnya berasal dari minyak bumi, yang
tersimpan di bawah tanah selama jutaan tahun: membakarnya menambah
karbon dioksida ke \omterm{def:g4:air-a-mixture-of-gases:air}{udara}, seperti membakar \omterm{def:g8:combustion-and-fuels:fossil}{bahan bakar fosil}.
\end{solution}

\begin{solution}{exo:g8:plastics:12}
Botol itu tidak membusuk: matahari dan ombak memecahnya menjadi
kepingan yang makin lama makin kecil. Kepingan-kepingan kecil itu
terapung atau tenggelam di dalam air, dan ikan menelannya bersama
makanannya.
\end{solution}

\begin{solution}{pb:g8:plastics:1}
\textbf{1.} Badan botol: 1 (PET). Tutup: 2 (HDPE) atau 5 (PP).

\textbf{2.} Sintetis: \omterm{def:g8:plastics:plastic}{plastik} itu dibuat dari zat-zat yang dibangun para
ahli kimia, sebagian besar dari minyak bumi.

\textbf{3.} Ya: PET adalah \omterm{def:g8:plastics:thermoplastic}{termoplastik}.

\textbf{4.} Setiap \omterm{def:g8:plastics:plastic}{plastik} didaur ulang tersendiri: jika dilelehkan
bersama, \omterm{def:g6:pure-substances-and-mixtures:mixture}{campuran} \omterm{def:g8:plastics:plastic}{plastik} menghasilkan bahan yang buruk.

\textbf{5.} $0.80 \times 1000 = \qty{800}{kg}$.

\textbf{6.} Tutup botol: $0.12 \times 1000 = \qty{120}{kg}$. Label: $0.05
\times 1000 = \qty{50}{kg}$.

\textbf{7.} $0.03 \times 1000 = \qty{30}{kg}$ kotoran dan cairan.

\textbf{8.} $80 + 12 + 5 = \qty{97}{\percent}$.

\textbf{9.} $0.025 \times 800 = \qty{20}{kg}$.

\textbf{10.} Serpihan itu menghemat minyak bumi dan energi untuk membuat
PET baru, dan mengubah sampah menjadi benda yang berguna.

\textbf{11.} Karbon dioksida (dan, jika \omterm{def:g4:what-a-fire-needs:burning}{pembakarannya} buruk, karbon
monoksida dan jelaga).

\textbf{12.} $800 - 20 = \qty{780}{kg}$ serpihan PET bersih per bongkah.
\end{solution}
